\chapter{שיטת ה-\texorpdfstring{\en{Lean Startup}}{Lean Startup} ולולאת \texorpdfstring{\en{Build--Measure--Learn}}{Build-Measure-Learn}}

\emph{פרק זה מדגים מעבר דו-כיווני} \term{(BiDi)} \emph{נכון בין עברית מימין-לשמאל
לבין מונחי אנגלית וקוד משמאל-לימין.}

שיטת ה-\en{Lean Startup} של אריק ריס מציעה מנוע צמיחה המבוסס על
\emph{למידה מאומתת} \term{(validated learning)}: כל מחזור פיתוח מתחיל בהשערה,
ממיר אותה למוצר מינימלי בר-קיימא — \term{Minimum Viable Product (MVP)} —
ומודד את תגובת הלקוחות בפועל~\cite{ries2011lean}.

\section{לולאת \texorpdfstring{\en{Build--Measure--Learn}}{Build-Measure-Learn}}

ליבת השיטה היא לולאת משוב מהירה: \en{Ideas} $\rightarrow$ \en{Build}
$\rightarrow$ \en{Product} $\rightarrow$ \en{Measure} $\rightarrow$ \en{Data}
$\rightarrow$ \en{Learn}, וחוזר חלילה. המטרה אינה לעבור את הלולאה פעם אחת, אלא
\emph{למזער את הזמן} למחזור שלם. איור~\ref{fig:bml} מתאר את הלולאה.

\begin{figure}[h]
  \centering
  \begin{tikzpicture}[every node/.style={font=\small\itshape}]
    \node[draw, circle, minimum size=1.5cm, fill=brandblue!10] (build) at (0,0) {\en{Build}};
    \node[draw, circle, minimum size=1.5cm, fill=brandgreen!10] (measure) at (3.4,-2.2) {\en{Measure}};
    \node[draw, circle, minimum size=1.5cm, fill=brandred!10] (learn) at (-3.4,-2.2) {\en{Learn}};
    \draw[-{Stealth[length=3mm]}, very thick, brandblue]
      (build) to[bend left=20] node[right=2pt]{\en{Product}} (measure);
    \draw[-{Stealth[length=3mm]}, very thick, brandgreen]
      (measure) to[bend left=20] node[below=2pt]{\en{Data}} (learn);
    \draw[-{Stealth[length=3mm]}, very thick, brandred]
      (learn) to[bend left=20] node[left=2pt]{\en{Ideas}} (build);
  \end{tikzpicture}
  \caption{לולאת \en{Build--Measure--Learn} — מנוע הלמידה של ה-\en{Lean Startup},
  משורטטת ב-\en{TikZ}.}
  \label{fig:bml}
\end{figure}

\section{\texorpdfstring{\en{MVP}}{MVP}, פיבוט ומדדים שאינם משקרים}

ה-\en{MVP} הוא הגרסה הזולה ביותר שעדיין מאפשרת לבדוק את ההשערה המרכזית — לא
מוצר חצי-אפוי, אלא \emph{ניסוי ממוקד}. כאשר הנתונים מלמדים שההשערה שגויה, מבצעים
\emph{פיבוט} \term{(pivot)} — שינוי מבני בכיוון תוך שמירה על מה שנלמד. חשוב להבחין
בין \term{actionable metrics} (מדדים מובילים לפעולה) לבין \term{vanity metrics}
(מדדי יוהרה) כגון מספר הורדות מצטבר, שנראים טוב אך אינם מנחים החלטה.

לדוגמה, השערה נבדקת מנוסחת באנגלית כשורה אחת:
\begin{quote}\en{\codefont If we add one-click checkout, weekly repeat purchases rise by 15\%.}\end{quote}
לאחר ההגדרה הזו חוזרים לעברית: רק כאשר המספר מאומת בנתוני אמת, מקדמים את הפיצ׳ר
משלב הניסוי אל המוצר.
